\section{导读：这期季报的主线是什么}

本节先建立整期的阅读框架。张小珺的“全球大模型季报”不是逐条新闻播报，而是用一组判断把季度变化压缩成技术路线图。2025 年 Q1 的关键变化，是李广密重新强调 Pre-training 的基础作用，同时把 Coding、Agent、Online Learning、模型公司战略分化和中美格局放在同一条 AGI 主线上讨论。

\begin{importantbox}{本期核心命题}
广密的核心判断可以压缩成一句话：智能提升仍然是唯一主线，Pre-training 负责打开模型内在上限，Coding 提供最通用的数字行动环境，Agent 是新物种，Online Learning 可能是下一条范式级路线；产品公司和模型公司都必须围绕这些能力重排自己的位置。
\end{importantbox}

\begin{knowledgebox}{视觉策略说明}
这期视频是固定播客画面，没有投屏、白板或产品演示。正文只把封面用于来源识别；正文图像全部是自制概念图，用来解释技术路线、组织战略、产品壁垒和产业格局。这样比重复主持人/嘉宾画面更符合本仓库的播客处理标准。
\end{knowledgebox}

\begin{figure}[H]
\centering
\includegraphics[width=0.92\textwidth]{q1-model-map.png}
\caption{2025 Q1 大模型季报地图：Pre-training、Coding、Agent、Online Learning 与全球格局共同构成主线。自制概念图，依据 00:00:08--00:04:22 与全片内容整理。}
\end{figure}

\begin{knowledgebox}{读图：不要把季报读成新闻清单}
图中 Q1 2025 位于中心，周围是 Pre-training、Coding、Agent、Online、产品和格局。读这张图时先看依赖关系：模型能力决定可用上限，Coding 和 Agent 提供行动环境，产品承接智能红利，公司格局则反映组织如何押注这些路线。
\end{knowledgebox}

\subsection{阅读路线}

本节给出阅读路线。全片可以拆成六个问题：第一，为什么在大家转向应用和后训练时，广密重新强调 Pre-training？第二，为什么 Coding 不只是编程，而是 AGI 早期最重要的环境？第三，OpenAI 和 Anthropic 的战略分化说明了什么组织问题？第四，Agent 为什么被称为新物种？第五，Online Learning 是否可能改变训练范式？第六，模型、产品、资本和中美格局如何互相约束？

\noindent\begin{tabularx}{\textwidth}{p{0.19\textwidth}p{0.34\textwidth}X}
\toprule
阅读问题 & 访谈中的材料 & 需要形成的判断 \\
\midrule
模型上限来自哪里 & Pre-training、Post-training、RL、合成数据 & 后训练可以塑形，但底座上限仍来自预训练与数据/算力/架构。 \\
Agent 如何落地 & Coding、computer use、tool use、long context & 数字环境提供可验证反馈，是 Agent 先落地的地方。 \\
公司为何分化 & OpenAI、Anthropic、DeepSeek、Manus、Cursor & 战略不是口号，而是组织能力、产品入口和技术积累的表达。 \\
产品壁垒在哪里 & Cloud、OS、容器、token 成本、盗火者 & 裸模型时代变弱，承接智能的环境和工作流变强。 \\
下一范式是什么 & Online Learning、long-term memory、环境反馈 & 如果模型能边行动边学习，智能提升路径会再次改变。 \\
\bottomrule
\end{tabularx}

\begin{warningbox}{重要边界：这是观点型季报，不是事实年鉴}
本笔记把嘉宾观点整理成分析框架，但不把“AGI 两年内实现”“某公司组合权重”“某公司会失败”这类判断写成事实。另一个术语边界是：本期出现的 Perplexity 指 AI 搜索产品/公司，不是语言模型评测里的 perplexity（PPL，交叉熵取指数后得到的困惑度指标，可直觉理解为每个 token 的有效候选数）。
\end{warningbox}

\subsection{本章小结}

这期季报的技术价值，在于它把 2025 年 Q1 的几条热门线索重新组织为一条 AGI 主线：模型上限、数字行动环境、Agent 产品化、在线学习范式、公司战略和地缘约束。后文会逐层展开。

